% ============================================================
%  part7/ch36-cluster-formation.tex
%  Chapter 36: Lensing, Void Statistics, Filament Statistics
% ============================================================

\chapter{Lensing, Void Statistics, and Filament Statistics}
\label{ch:cluster-formation}

\begin{WhyChapter}
  Gravitational lensing, void statistics, and filament density
  are three of the most directly measurable large-scale structure
  observables.
  This chapter derives RSVP predictions for each and compares
  them with $\Lambda$CDM.
  The RSVP interpretations are significantly different in their
  causal mechanism while remaining observationally competitive
  at current precision.
\end{WhyChapter}

\section{Gravitational lensing from admissibility gradients}

Define effective refractive index
\[
  n(x) = 1 + \epsilon\Phi(x) + \eta L(x).
\]
Light follows $\delta\int n\,ds = 0$.

\begin{theorem}[RSVP lensing deflection]
  \label{thm:lensing}
  \statusproven{}
  The deflection angle satisfies
  \[
    \theta
    \approx \epsilon\int\nabla_\perp\Phi\,ds
    + \eta\int\nabla_\perp L\,ds.
  \]
\end{theorem}

\begin{proof}
  The ray equation $\frac{d}{ds}(n\dot{x}) = \nabla n$ gives,
  for weak fields,
  $\theta = \int\nabla_\perp n\,ds$.
  Substituting the refractive index definition gives the result.
\end{proof}

\begin{WorkedExample}[Lamphron contribution to lensing]
  Let $\int\nabla_\perp\Phi\,ds = 2$,
  $\int\nabla_\perp L\,ds = 5$,
  $\epsilon = 1$, $\eta = 0.5$.
  Then $\theta = 2 + 2.5 = 4.5$.

  More than half the deflection originates from lamphron structure.
  \textbf{Interpretation:} An observer measuring lensing cannot
  determine directly whether deflection arose from visible matter
  or hidden plenum structure.
  This is a reconstruction problem, not a measurement problem.
\end{WorkedExample}

\begin{ReconView}
  Observable: deflection angle $\theta$.
  Hidden: $(\Phi, L, S, \mathbf{v})$.
  Measurement: $\pi(F(\Phi, L)) = \theta$.
  Inverse problem: find admissible field configurations
  $(\Phi, L)$ consistent with the observed $\theta$.
  Multiple configurations may produce the same deflection.
\end{ReconView}

\begin{Falsifiability}
  If cross-correlation between lensing maps and spectroscopic
  surveys shows that $\theta$ traces the matter distribution
  with no residual ($\eta = 0$), the lamphron contribution
  to lensing would be excluded.
\end{Falsifiability}

\section{Void statistics and the coarsening exponent}

\begin{theorem}[Void size distribution scaling]
  \label{thm:void-scaling}
  \statusconjectured{}
  Under RSVP coarsening, void radii satisfy dynamic scaling:
  \[
    N(R,t) = R^{-\tau}\,f\!\left(\frac{R}{t^{1/z}}\right).
  \]
\end{theorem}

\begin{proof}[Proof sketch]
  Standard coarsening theory implies self-similarity with
  characteristic length $\ell(t) \sim t^{1/z}$.
  Identifying void domains with coarsening domains gives the
  stated distribution.
\end{proof}

\begin{WorkedExample}[Void growth under $z=2$]
  For $z = 2$: $\ell(t) \sim t^{1/2}$.
  A void of characteristic radius 10 at $t = 100$ has
  radius $\approx 20$ at $t = 400$.

  \textbf{Interpretation:} Void radii grow as $\sqrt{t}$ under
  Allen--Cahn-dominated dynamics.
  Void catalogues at multiple redshifts can estimate
  $z_{\mathrm{RSVP}}$ directly from the growth rate.
\end{WorkedExample}

\begin{WorkedExample}[Empirical determination of $z$]
  Suppose two surveys yield $\rho_f(t_1) = 100$ and
  $\rho_f(t_2) = 50$.
  Then $\rho_f(t_2)/\rho_f(t_1) = (t_2/t_1)^{-2/z}$,
  giving $z = -2\log(t_2/t_1)/\log(1/2)$.
  This is a direct empirical estimate of $z_{\mathrm{RSVP}}$.
\end{WorkedExample}

\section{Filament statistics}

\begin{proposition}[Filament density scaling]
  \label{prop:filament-density}
  \statusconjectured{}
  Filament (defect) density scales as
  \[
    \rho_f \sim \ell^{-2} \sim t^{-2/z}.
  \]
\end{proposition}

\begin{proof}[Proof sketch]
  Defects occupy codimension two.
  Their number per unit volume scales inversely with domain area
  $\ell^2$.
\end{proof}

\begin{WorkedExample}[Filament decay for $z=2$]
  For $z = 2$: $\rho_f \sim t^{-1}$.
  Doubling cosmic age halves filament density.
\end{WorkedExample}

\begin{NotUnique}
  The filament density scaling $\rho_f \sim t^{-2/z}$ is
  mathematically identical to vortex annihilation in superfluids
  and domain-wall decay in Ising systems.
  The interpretation changes; the scaling law does not.
  RSVP does not invent new mathematics.
  It applies established coarsening universality classes to
  cosmological structure.
\end{NotUnique}

\begin{ChapterConnection}
  The filament statistics derived here are the cosmological
  analogue of the defect-density results from
  Chapter~\ref{ch:defects}.
  The exponent $z$ is the same object computed in
  Chapter~\ref{ch:cahn-hilliard} for Cahn--Hilliard dynamics
  and conjectured to differ in Chapter~\ref{ch:entropy-production}
  for the full RSVP system.
  The cosmic web becomes a single member of the universal family
  of coarsening defect networks.
\end{ChapterConnection}

\begin{ProblemSet}
\problem{
  For RSVP lensing with $\epsilon = 1$ and $\eta = 0.5$,
  compute the deflection angle if $\int\nabla_\perp\Phi\,ds = 3$
  and $\int\nabla_\perp L\,ds = 0$.
  Interpret: what does zero lamphron gradient mean observationally?
}
\problem{
  Given the void size distribution
  $N(R,t) = R^{-\tau} f(R/t^{1/z})$,
  show that the mean void radius $\langle R\rangle(t) \sim t^{1/z}$.
  What does this imply for the time-evolution of the void
  size function?
}
\problem{
  Using the empirical determination method from
  Worked Example~\ref{ch:cluster-formation}.2,
  estimate $z_{\mathrm{RSVP}}$ if filament densities are
  $\rho_f = 80$ at $t_1$ and $\rho_f = 20$ at $t_2 = 4t_1$.
}
\problem{
  \textbf{Inverse problem.}
  Suppose a weak lensing survey measures $\theta = 3.5$
  along a line of sight.
  If $\epsilon = 1$ and the scalar contribution is estimated
  as $\epsilon\int\nabla_\perp\Phi\,ds = 2.0$,
  what is the admissible range for the lamphron contribution
  $\int\nabla_\perp L\,ds$?
  What additional observations would narrow this range?
}
\end{ProblemSet}
